\section{USAMO 1988}
        \begin{enumerate}
        	\item (\textit{USAMO 1988}) The repeating decimal $0.ab\cdots k\overline{pq\cdots u}=\frac mn$, where $m$ and $n$ are relatively prime integers, and there is at least one decimal before the repeating part. Show that $n$ is divisible by 2 or 5 (or both). (For example, $0.011\overline{36}=0.01136363636\cdots=\frac 1{88}$, and 88 is divisible by 2.)


 



\textbf{Solution 1}

First, split up the nonrepeating parts and the repeating parts of the decimal, so that the nonrepeating parts equal to $\frac{a}{b}$ and the repeating parts of the decimal is equal to $\frac{c}{d}$.

 Suppose that the length of  $0.ab\cdots k$ is $p$ digits. Then $\frac{a}{b} = \frac{0.ab\cdots k}{10^{p+1}}$ Since $0.ab\cdots k < 10^{p+1}$, after reducing the fraction, there MUST be either a factor of 2 or 5 remaining in the denominator. After adding the fractions $\frac{a}{b}+\frac{c}{d}$, the simplified denominator $n$ will be $\\lcm(b,d)$ and since $b$ has a factor of $2$ or $5$,  $n$ must also have a factor of 2 or 5.

 $\blacksquare$

 



\textbf{Solution 2}

It is well-known that $0.ab...k \overline{pq...u} = \frac{ab...u - ab...k}{99...900...0}$, where there are a number of 9s equal to the count of digits in $pq...u$, and there are a number of 0s equal to the count of digits $c$ in $ab...k$. Obviously $ab...k$ is different from $pq...u$ (which is itself the repeating part), so the numerator cannot have $c$ consecutive terminating zeros, and hence the denominator still possesses a factor of 2 or 5 not canceled out. This completes the proof.

 


	\item (\textit{USAMO 1988}) The cubic polynomial $x^3+ax^2+bx+c$ has real coefficients and three real roots $r\ge s\ge t$. Show that $k=a^2-3b\ge 0$ and that $\sqrt k\le r-t$.


 



\textbf{Solution 1}

By Vieta's Formulas, $a=-r-s-t$, $b=rs+st+rt$, and $c=-rst$.Now we know $k=a^2-3b$; in terms of r, s, and t, then, 
 \[k=(-r-s-t)^2-3(rs+st+rt)\]
 \[k=r^2+s^2+t^2-rs-st-rt\]Now notice that we can multiply both sides by 2, and rearrange terms to get$2k=(r-s)^2+(s-t)^2+(r-t)^2$.But since $r, s, t\in \mathbb{R}$, the three terms of the RHS are all non-negative (as the square of a real number is always non-negative), and therefore their sum is also non-negative -- that is, $2k\ge 0 \Rightarrow k\ge 0$.

 Now, we will show that $\sqrt k\le r-t$.We can square both sides, and the inequality will hold since they are both non-negative (it is given that $r\ge t$, therefore $r-t\ge 0$). This gives $k \le r^2-2rt+t^2$.Now we already have $k=r^2+s^2+t^2-rs-st-rt$, so substituting this for k gives
 \[r^2+s^2+t^2-rs-st-rt \le r^2-2rt+t^2\]
 \[s^2-rs-st+rt \le 0\]
 \[s^2-(r+t)s+rt \le 0\]Note that this is a quadratic. Since its leading coefficient is positive, its value is less than 0 when s is between the two roots. Using the quadratic formula:
 \[s=\frac {r+t\pm \sqrt {(-r-t)^2-4rt} } 2\]
 \[s=\frac {r+t\pm \sqrt {r^2-2rt+t^2} } 2\]
 \[s=\frac {r+t\pm (r-t) } 2\]
 \[s \in \{r, t\}\]The quadratic is 0 when s is equal to r or t, and the inequalityholds when its value is less than or equal to 0 -- that is, $r\ge s\ge t$.(Its value is less than or equal to 0 when s is between the roots, since thegraph of the quadratic opens upward.)In fact, the problem tells us this is true. Q.E.D.

 



\textbf{Solution 2}

From Vieta's Formula (which tells us that $a = -(r+s+t)$ and $b = rs + st + rt$), we have that
 \[k = a^2 - 3b = r^2 + s^2 + t^2 - rs - st - rt = \frac{1}{2} ((r-s)^2 + (s-t)^2 + (r-t)^2),\]clearly non-negative. To prove $\sqrt{k} \le r - t$, it suffices to prove the square of this relation, or 
 \[r^2 + s^2 + t^2 - rs - st - rt \le r^2 - 2rt + t^2.\] This in turn simplifies to 
 \[rs + st - rt - s^2 \ge 0,\] or 
 \[(r - s)(s - t) \ge 0,\] which is clearly true as $r \ge s \ge t$. This completes the proof.

 



\textbf{Solution 3}

By Vieta's Formulas, $a = -(r+s+t)$ and $b = rs + st + rt$. $k = (r+s+t)^2 - 3(rs + st + rt) = r^2+s^2+t^2-rs-st-rt = \frac{(r-s)^2 + (s-t)^2 + (r-t)^2}{2}$.

 To show that $k \ge 0$, simply note that by the trivial inequality, all three squares are greater than $0$ as they are the squares of real numbers.

 To show that $\sqrt{k} \le r-t$, since both are positive, it is sufficient to show that $k \le (r-t)^2$. $\frac{(r-s)^2 + (s-t)^2 + (r-t)^2}{2} \le (r-t)^2$ implies that $k \le (r-t)^2$. $\frac{(r-s)^2 + (s-t)^2 - (r-t)^2}{2} \le 0$. Let $y = r-s$ and $z = s-t$. We then have $\frac{y^2 + z^2 - (y+z)^2}{2} \le 0 \implies -2yz \le 0$, which is clearly true as both $y$ and $z$ are positive.

 


	\item (\textit{USAMO 1988}) Let $X$ be the set $\{ 1, 2, \cdots , 20\}$ and let $P$ be the set of all 9-element subsets of $X$. Show that for any map $f: P\mapsto X$ we can find a 10-element subset $Y$ of $X$, such that $f(Y-\{k\})\neq k$ for any $k$ in $Y$.


 



\textbf{Solution}

For a given function $f$, consider the set of elements $Q \in P$ paired with elements $t \in S$ such that $f(Q) \ne t$. Call the set of pairs $(Q,t)$ set $R$.

 To employ the Pigeonhole Principle, we need a lower bound on $\lvert R \rvert$. There are $\binom{20}{9}$ elements $Q$ in $P$. If $f(Q) \in Q$, then there are $11$ choices for $t$ so that $f(Q) \ne t$. If $f(Q) \notin Q$, then there are $10$ choices for $t$ so that $f(Q) \ne t$. Therefore $\lvert R \rvert \ge 10 \cdot \binom{20}{9}$. These pairs $(Q,t)$ must be dispersed over all the $10$-element subsets of $S$. By the Pigeonhole Principle, some $10$-element subset of $S$ must contain at least 

 


	\item (\textit{USAMO 1988}) $\Delta ABC$ is a triangle with incenter $I$. Show that the circumcenters of $\Delta IAB$, $\Delta IBC$, and $\Delta ICA$ lie on a circle whose center is the circumcenter of $\Delta ABC$.


 



\textbf{Solution 1}

Let the circumcenters of $\Delta IAB$, $\Delta IBC$, and $\Delta ICA$ be $O_c$, $O_a$, and $O_b$, respectively. It then suffices to show that $A$, $B$, $C$, $O_a$, $O_b$, and $O_c$ are concyclic.

 We shall prove that quadrilateral $ABO_aC$ is cyclic first. Let $\angle BAC=\alpha$, $\angle CBA=\beta$, and $\angle ACB=\gamma$. Then $\angle ICB=\gamma/2$ and $\angle IBC=\beta/2$. Therefore minor arc $\overarc{BIC}$ in the circumcircle of $IBC$ has a degree measure of $\beta+\gamma$. This shows that $\angle CO_aB=\beta+\gamma$, implying that $\angle BAC+\angle BO_aC=\alpha+\beta+\gamma=180^{\circ}$. Therefore quadrilateral $ABO_aC$ is cyclic.

 This shows that point $O_a$ is on the circumcircle of $\Delta ABC$. Analagous proofs show that $O_b$ and $O_c$ are also on the circumcircle of $ABC$, which completes the proof. $\blacksquare$

 



\textbf{Solution 2}

Let $M$ denote the midpoint of arc $AC$. It is well known that $M$ is equidistant from $A$, $C$, and $I$ (to check, prove $\angle IAM = \angle AIM = \frac{\angle BAC + \angle ABC}{2}$), so that $M$ is the circumcenter of $AIC$. Similar results hold for $BIC$ and $CIA$, and hence $O_c$, $O_a$, and $O_b$ all lie on the circumcircle of $ABC$.

 



\textbf{Solution 3}

Extend $CI$ to point $L$ on $(ABC)$. By The Incenter-Excenter Lemma, B, I, A are all concyclic. Thus, L is the circumcenter of triangle $IAB$. In other words, $L=O_c$, so $O_c$ is on $(ABC)$. Similarly, we can show that $O_a$ and $O_b$ are on $(ABC)$, and thus, $A,B,C,O_a,O_b,O_c$ are all concyclic. It follows that the circumcenters are equal.

 



\textbf{Solution 4}

Let the centers be $T, R, S$. We want to show that $ABC$ and $TRS$ have the same circumcircle. By Fact 5 we know that $CATB$ lie on a circle and similarly with the others. Thus the two triangles have the same circumcircle.~coolmath\_2018 

 


	\item (\textit{USAMO 1988}) Let $p(x)$ be the polynomial $(1-x)^a(1-x^2)^b(1-x^3)^c\cdots(1-x^{32})^k$, where $a, b, \cdots, k$ are integers. When expanded in powers of $x$, the coefficient of $x^1$ is $-2$ and the coefficients of $x^2$, $x^3$, ..., $x^{32}$ are all zero. Find $k$.


 



\textbf{Solution 1}

First, note that if we reverse the order of the coefficients of each factor, then we will obtain a polynomial whose coefficients are exactly the coefficients of $p(x)$ in reverse order. Therefore, if
 \[p(x)=(1-x)^{a_1}(1-x^2)^{a_2}(1-x^3)^{a_3}\cdots(1-x^{32})^{a_{32}},\]we define the polynomial $q(x)$ to be
 \[q(x)=(x-1)^{a_1}(x^2-1)^{a_2}(x^3-1)^{a_3}\cdots(x^{32}-1)^{a_{32}},\]noting that if the polynomial has degree $n$, then the coefficient of $x^{n-1}$ is $-2$, while the coefficients of $x^{n-k}$ for $k=2,3,\dots, 32$ are all $0$.

 Let $P_n$ be the sum of the $n$th powers of the roots of $q(x)$. In particular, by Vieta's formulas, we know that $P_1=2$. Also, by Newton's Sums, as the coefficients of $x^{n-k}$ for $k=2,3,\dots,32$ are all $0$, we find that
 \begin{align*} P_2-2P_1&=0\\ P_3-2P_2&=0\\ P_4-2P_3&=0\\ &\vdots\\ P_{32}-2P_{31}&=0. \end{align*}
Thus $P_n=2^n$ for $n=1,2,\dots, 32$. Now we compute $P_{32}$. Note that the roots of $(x^n-1)^{a_n}$ are all $n$th roots of unity. If $\omega=e^{2\pi i/n}$, then the sum of $32$nd powers of these roots will be
 \[a_n(1+\omega^{32}+\omega^{32\cdot 2}+\cdots+\omega^{32\cdot(n-1)}).\]If $\omega^{32}\ne 1$, then we can multiply by $(\omega^{32}-1)/(\omega^{32}-1)$ to obtain
 \[\frac{a_n(1-\omega^{32n})}{1-\omega^{32}}.\]But as $\omega^n=1$, this is just $0$. Therefore the sum of the $32$nd powers of the roots of $q(x)$ is the same as the sum of the $32$nd powers of the roots of
 \[(x-1)^{a_1}(x^2-1)^{a_2}(x^4-1)^{a_4}(x^{8}-1)^{a_4}(x^{16}-1)^{a_{16}}(x^{32}-1)^{a_{32}}.\]The $32$nd power of each of these roots is just $1$, hence the sum of the $32$nd powers of the roots is
 \[P_{32}=2^{32}=a_1+2a_2+4a_4+8a_8+16a_{16}+32a_{32}.\tag{1}\]On the other hand, we can use the same logic to show that
 \[P_{16}=2^{16}=a_1+2a_2+4a_4+8a_8+16a_{16}.\tag{2}\]Subtracting (2) from (1) and dividing by 32, we find
 \[a_{32}=\frac{2^{32}-2^{16}}{2^5}.\]Therefore, $a_{32}=2^{27}-2^{11}$.

 


 



\textbf{Solution 2}

By a limiting process, we can extend the problem to that of finding a sequence $b_1, b_2, \ldots$ of integers such that
 \[(1 - z)^{b_1}(1 - z^2)^{b_2}(1 - z^3)^{b_3}\cdots = 1 - 2z.\](The notation comes from the Alcumus version of this problem.)

 If we take logarithmic derivatives on both sides, we get
 \[\sum_{n = 1}^{\infty}\frac{b_n\cdot (-nz^{n - 1})}{1 - z^n} = \frac{-2}{1 - 2z},\]and upon multiplying both sides by $-z$, this gives us the somewhat simple form
 \[\sum_{n = 1}^{\infty} nb_n\cdot\frac{z^n}{1 - z^n} = \frac{2z}{1 - 2z}.\]Expanding all the fractions as geometric series, we get
 \[\sum_{n = 1}^{\infty} nb_n\sum_{k = 1}^{\infty} z^{nk} = \sum_{n = 1}^{\infty} 2^nz^n.\]Comparing coefficients, we get
 \[\sum_{d\mid n} db_d = 2^n\]for all positive integers $n$. In particular, as in Solution 1, we get
 \[\begin{array}{ll} b_1 + 2b_2 + 4b_4 + 8b_8 + 16b_{16} + 32b_{32} &= 2^{32}, \\ b_1 + 2b_2 + 4b_4 + 8b_8 + 16b_{16}\phantom{ + 32b_{32}} &= 2^{16}, \end{array}\]from which the answer $b_{32} = 2^{27} - 2^{11}$ follows.

 \textbf{Remark:} To avoid the question of what an infinite product means in the context of formal power series, we could instead view the problem statement as saying that
 \[(1 - z)^{b_1}(1 - z^2)^{b_2}\cdots (1 - z^{32})^{b_{32}}\equiv 1 - 2z\pmod{z^{33}};\]modular arithmetic for polynomials can be defined in exactly the same way as modular arithmetic for integers. Uniqueness of the $b_n$'s comes from the fact that we have
 \[(1 - z)^{b_1}\cdots (1 - z^{n - 1})^{b_{n - 1}}\equiv 1 - 2z\pmod{z^n}\]for all $n\leq 33$ by further reduction modulo $z^n$ (as $z^n\mid z^{33}$ for $n\leq 33$), so we could uniquely solve for the $b_n$'s one at a time. (This idea can be pushed further to explain why it's fine to pass to the infinite product version of the problem.)

 To convert the above solution to one that works with polynomials modulo $z^{33}$, note that the derivative is not well-defined, as for instance, $1$ and $1 + z^{33}$ are equivalent modulo $z^{33}$, but their derivatives, $0$ and $33z^{32}$, are not. However, the operator $f(z)\mapsto zf'(z)$ \textit{is} well-defined. The other key idea is that for any $n$, modulo $z^n$, polynomials of the form $1 - zf(z)$ are invertible, with inverse
 \[\frac{1}{1 - zf(z)}\equiv\frac{1 - (zf(z))^n}{1 - zf(z)} = 1 + zf(z) + \cdots + (zf(z))^{n - 1}).\]Therefore, for the polynomial in the problem, call it $g(z)$, we can still form the expression $zg'(z)/g(z)$, which is what we originally got by taking the logarithmic derivative and multiplying by $z$, and expand it to eventually get
 \[\sum_{n = 1}^{32} nb_n\sum_{k = 1}^{32} z^{nk}\equiv\sum_{n = 1}^{32} 2^nz^n\pmod{z^{33}},\]which gets us the same relations (for $n\leq 32$).

 



\textbf{Solution 3}

From the starting point of Solution 2,
 \[(1 - z)^{b_1}(1 - z^2)^{b_2}(1 - z^3)^{b_3}\cdots = 1 - 2z,\]taking reciprocals and expanding with geometric series gives us
 \[\prod_{n = 1}^{\infty}\left(\sum_{k = 0}^{\infty} z^{kn}\right)^{b_n} = \sum_{n = 0}^{\infty} 2^nz^n.\]On the right, we have the generating function for the number of monic polynomials of degree $n$ over the field $\mathbb{F}_2$ of two elements, and on the left, we have the factorisation of this generating function that considers the breakdown of any given monic polynomial into monic irreducible factors. As such, we have the interpretation
 \[b_n = \text{number of monic irreducible polynomials of degree }n\text{ over }\mathbb{F}_2.\]From here, to determine $b_n$, we analyse the elements of $\mathbb{F}_{2^n}$, of which there are $2^{n}$ in total. Given $\alpha\in\mathbb{F}_{2^n}$, if the minimal polynomial $f_{\alpha}$ of $\alpha$ has degree $d$, then $d\mid n$ and all other roots of $f_{\alpha}$ appear in $\mathbb{F}_{2^n}$. Moreover, if $d\mid n$ and $f$ is an irreducible polynomial of degree $d$, then all roots of $f$ appear in $\mathbb{F}_{2^n}$. (These statements are all well-known in the theory of finite fields.) As such, for each $d\mid n$, there are precisely $db_d$ elements of $\mathbb{F}_{2^n}$ of degree $d$, and we obtain the same equation as in Solution 2,
 \[\sum_{d\mid n} db_d = 2^n.\]The rest is as before.

 


 


\end{enumerate}